%%=============================================================================
%% CAPITULO 5 -- CONSIDERACOES FINAIS
%%=============================================================================

\section{Considerações finais}

Este trabalho perguntou qual estratégia de recuperação --- RAG denso, RAG
denso-esparso ou GraphRAG incremental --- e qual estratégia de geração ---
single-pass ou verificada multiagente --- produz respostas mais precisas,
fundamentadas e rastreáveis sobre um corpus jurídico-administrativo
brasileiro. A resposta obtida diverge, em parte, da hipótese que orientou o
desenho do estudo: o GraphRAG incremental não supera o RAG denso-esparso com
significância estatística em nenhuma das quatro métricas avaliadas, apesar
de superar o RAG denso puro nas quatro; a geração verificada multiagente não
melhora a fidelidade nem a taxa de verificação de citação em relação à
geração single-pass, em nenhuma das duas bases de retrieval testadas.

O objetivo geral --- comparar empiricamente estratégias de recuperação e de
geração sobre esse corpus --- foi perseguido nesta fase do projeto através
do desenho descrito no Capítulo 3. Dos objetivos específicos, construir o
corpus e o harness de avaliação é uma frente já em andamento; avaliar H1 e
H2 depende da execução completa do harness sobre o golden dataset, ainda em
curadoria nesta etapa do mestrado.

A formatação deste trabalho segue as normas da ABNT para apresentação de
trabalhos acadêmicos \parencite{abnt14724}, para elaboração de referências
\parencite{abnt6023} e para apresentação do sumário \parencite{abnt6027}. Note,
na lista de referências, que as entradas repetidas do mesmo autor são
substituídas por um travessão, conforme exige a NBR 6023.
